% !TeX program = xelatex
% ==========================================================================
% EXAMPLE of a filled-in proposal. Fictional project, fictional students.
% Demonstrates the expected level of detail: a concrete problem, sources,
% named datasets, concrete metrics and a baseline, a realistic plan.
% ==========================================================================
\documentclass[english]{ucuproposal}

\projecttitle{UA-Digest: Abstractive Summarization of Ukrainian News}
\projecttagline{Compact, factually accurate digests of Ukrainian news articles using a fine-tuned language model}

\teammember[https://www.linkedin.com/in/ivan-petrenko]{Ivan Petrenko}
\teammember{Olena Bondarenko}
\teammember{Andrii Kovalenko}

\mentor{Iryna Sydorenko}

\track{engineering}

\repo{} % no repository yet -- team hasn't set one up at this stage

\proposaldate{2026-09-21}

\begin{document}
\maketitlepage

\section{Problem and motivation}
There are almost no tools for automatic news summarization in Ukrainian:
available solutions either target English or use extractive methods that
produce disjointed text. Ukrainian media publish hundreds of articles a
day, and readers lack a quick way to get a compact, factually accurate
digest. Existing browser extensions for other languages rely on
extractive sentence selection, which tends to string together
grammatically disconnected fragments and loses the narrative flow that
readers expect from a lead paragraph. The project has practical value (a
plugin for reading news) and research value (evaluating the quality of
abstractive summarization for a low-resource, morphologically rich
language where labeled summarization data is scarce).

\section{Related work}
Abstractive summarization is based on seq2seq architectures with an
attention mechanism~\cite{vaswani2017attention}. BERT-based models and
their distilled versions~\cite{devlin2019bert,sanh2019distilbert} show
strong results for English but require adaptation for the morphologically
rich Ukrainian language: subword tokenizers trained on English corpora
fragment Ukrainian words far more aggressively, which hurts both training
efficiency and generation fluency. ROUGE~\cite{lin2004rouge} remains the
standard metric for evaluating summary quality, which we will supplement
with a manual expert assessment of factual accuracy, since ROUGE alone
does not penalize hallucinated facts that are lexically similar to the
source text.

\section{Data}
The primary dataset is an open corpus of Ukrainian literary and
journalistic prose~\cite{ukrlit2023}, supplemented with a manually
collected set of \textasciitilde{}8\,000 ``article--headline+lead'' pairs
from three public RSS feeds of Ukrainian news outlets (usage license to
be confirmed with rights holders; for training we will use only articles
under an open license or with permission for research use). We will
deduplicate near-identical articles (syndicated wire content republished
by multiple outlets) using MinHash on word shingles, and filter out
articles shorter than 200 words, since these rarely contain enough
content for a meaningful lead. Split: 80\,/\,10\,/\,10\% (train/val/test),
stratified by outlet to avoid a single source dominating any one split.

\section{Approach and methods}
The base model is mT5-small, fine-tuned on the collected corpus using
teacher forcing with a cross-entropy loss. We plan to compare two
configurations: (1) direct fine-tuning of mT5-small, (2) mT5-small with
an additional self-supervised adaptation stage on an unannotated
Ukrainian corpus before fine-tuning. We use the Hugging Face Transformers
library implementation as a base, adapting tokenization and the sampling
scheme for Ukrainian. Training will use mixed-precision on a single GPU
with gradient accumulation to reach an effective batch size of 32,
early stopping on validation ROUGE-L, and beam search (beam size 4) at
inference time to reduce repetition artifacts common in small
fine-tuned models.

\section{Evaluation}
Metrics: ROUGE-1/2/L~\cite{lin2004rouge} against reference leads.
Baseline: extractive LexRank and a plain first-two-sentences cutoff
("lead-2"). Additionally, a manual evaluation of 50 random summaries by
three independent assessors on a 1--5 scale (coherence, factual
accuracy). We expect to beat the lead-2 baseline by at least 15\% on
ROUGE-L and achieve an average factual accuracy score of
\(\geq 4/5\). We will also report inter-annotator agreement (Krippendorff's
alpha) for the manual scores, and inspect a sample of low-scoring outputs
by hand to categorize the dominant failure modes (e.g. hallucinated
numbers vs. dropped named entities).

\section{Plan and risks}
\begin{keyfacts}
\textbf{Key numbers:} \textasciitilde{}8\,000 article pairs \textperiodcentered\ 2 model configurations \textperiodcentered\ 6 weeks to midterm
\end{keyfacts}
\textbf{Before midterm (6 weeks):} data collection and cleaning, basic
mT5-small fine-tuning, first ROUGE run against the baseline.
\textbf{Before the final (6 more weeks):} second model configuration,
manual evaluation, a simple web deployment demo. \textbf{Risks:}
(1) not enough high-quality annotated pairs -- plan B: expand the corpus
with semi-automatic lead generation; (2) low summary quality from
mT5-small -- plan B: switch to a distilled but larger model (mT5-base)
if compute resources allow; (3) rights holders decline permission for a
subset of the RSS-collected articles -- plan B: fall back to the open
literary/journalistic corpus alone and treat the RSS set as an optional
held-out evaluation sample rather than training data.

\section{Contribution statement}
\begin{contributions}
  \contribution{Ivan Petrenko}{35\%}{data collection and preprocessing, EDA, manual annotation}
  \contribution{Olena Bondarenko}{40\%}{model fine-tuning, experiments, hyperparameter tuning}
  \contribution{Andrii Kovalenko}{25\%}{training infrastructure, web deployment demo, documentation}
\end{contributions}

\ucuprintbibliography

\end{document}
